% LTeX: enabled=false

\documentclass{article}
\usepackage[utf8]{inputenc}

\newcommand{\nirpdftitle}{Lie Theory Review}
\usepackage{import}
\inputfrom{..}{nir}

\pagestyle{contentpage}

\title{Lie Theory for the Impatient}
\author{Nir Elber}
\date{Fall 2026}
\rhead{\textit{LIE THEORY SPEEDRUN}}
\lhead{}

\setcounter{tocdepth}{2}

\begin{document}

\maketitle

\begin{abstract}
	This document collects a variety of definitions and results from Lie theory. To avoid distraction, all vector spaces are over $\CC$, and all manifolds are real.
\end{abstract}

\tableofcontents

\newpage

\section{Lie Algebras}

\subsection{Starting Out}
\begin{definition}[Lie algebra]
	A \textit{Lie algebra} is a vector space $\mf g$ equipped with a \textit{Lie bracket} $[-,-]\colon\mf g\times\mf g\to\mf g$ satisfying the following properties.
	\begin{listroman}
		\item Bilinear: the map $[-,-]$ is bilinear.
		\item Alternating: $[X,X]=0$ for all $X\in\mf g$.
		\item Jacobi identity: $[X,[Y,Z]]+[Y,[Z,X]]+[Z,[X,Y]]=0$ for all $X,Y,Z\in\mf g$.
	\end{listroman}
	A morphism of Lie algebras is a linear map preserving the bracket.
\end{definition}
\begin{example}
	For any (possibly non-commutative, non-unital) algebra $A$ over $\CC$, one can define the Lie bracket $[-,-]\colon A\times A\to A$ by
	\[[a,b]\coloneqq ab-ba.\]
	For example, applying this to $A=\op{End}(V)$ where $V$ is a vector space gives the Lie algebra $\mf{gl}(V)$. If further $V=\CC^n$, we abbreviate $\mf{gl}(V)$ to $\mf{gl}_n(\CC)$.
\end{example}
\begin{example}
	The subspace
	\[\mf{sl}_n(\CC)\coloneqq\{X\in\mf{gl}_n(\CC):\tr X=0\}\]
	is a Lie subalgebra.
\end{example}
\begin{example}
	The subspace of alternating matrices is a Lie subalgebra, called $\mf o_n(\CC)\subseteq\mf{gl}_n(\CC)$.
\end{example}
\begin{definition}[representation]
	A \textit{representation} of a Lie algebra $\mf g$ is a vector space $V$ along with a morphism $\rho\colon\mf g\to\mf{gl}(V)$. A morphism of representations is a morphism of the underlying vector spaces which preserves the $\mf g$-action.
\end{definition}
\begin{example}[adjoint]
	The adjoint representation of a Lie algebra $\mf g$ is the morphism $\op{ad}_\bullet\colon\mf g\to\mf{gl}(\mf g)$ defined by
	\[\op{ad}_X(Y)\coloneqq[X,Y].\]
	The fact that $\op{ad}_\bullet$ preserves the Lie bracket is equivalent to the Jacobi identity.
\end{example}
\begin{definition}[irreducible]
	A representation $\rho\colon\mf g\to\mf{gl}(V)$ is \textit{irreducible} if and only if $V\ne0$, and $V$ admits no nonzero proper subrepresentations. A representation is \textit{completely reducible} if and only if it is a direct sum of irreducible representations.
\end{definition}
\begin{example}
	Let $\mf b\subseteq\mf{gl}_2(\CC)$ be the Lie subalgebra of upper-triangular matrices. Then the representation $\mf b\to\mf{gl}_2(\CC)$ is not completely reducible.
\end{example}
\begin{lemma}[Schur] \label{lem:schur}
	Let $\rho_V\colon\mf g\to\mf{gl}(V)$ and $\rho_W\colon\mf g\to\mf{gl}(W)$ be representations, and let $\varphi\colon V\to W$ be a nonzero morphism of representations.
	\begin{listalph}
		\item If $V$ is irreducible, then $\varphi$ is injective.
		\item If $W$ is irreducible, then $\varphi$ is surjective.
		\item If both $V$ and $W$ are irreducible, then $\varphi$ is an isomorphism.
	\end{listalph}
\end{lemma}
\begin{proof}[Sketch]
	Here, (c) follows from (a) and (b). We only sketch (a): because $\varphi$ is nonzero, $\ker\varphi$ is a proper subrepresentation of $V$, so $\ker\varphi=0$ because $V$ is irreducible.
\end{proof}
\begin{example}
	Let the natural action of $\mf{sl}_2(\CC)$ on $\CC^2$ is irreducible. It turns out that the irreducible representations of $\mf{sl}_2(\CC)$ are exactly the symmetric powers of $\CC^2$. To see this, show that $\op{diag}(+1,-1)$ acts semisimply on any irreducible representation, and then one can parameterize the irreducible representations by the highest eigenvalue of $\op{diag}(+1,-1)$.
\end{example}
\begin{definition}[universal enveloping algebra]
	Fix a Lie algebra $\mf g$. Then the \textit{universal enveloping algebra} $U\mf g$ of $\mf g$ is the free complex algebra generated by $\mf g$ along with the relations
	\[XY-YX=[X,Y]\]
	for each $X,Y\in\mf g$.
\end{definition}
\begin{remark}
	The data of a representation $\mf g\to\mathfrak{gl}(V)$ is equivalent to the data of a ring homomorphism $U\mf g\to\op{End}(V)$.
\end{remark}
\begin{theorem}[Poincar\'e--Birkhoff--Witt]
	Fix a Lie algebra $\mf g$, and let $\{X_i\}_{i=1}^n$ be an (ordered) basis. Then the monomials
	\[\{X_{i_1}\cdots X_{i_m}:i_1\le i_2\le\cdots\le i_m\}\]
	form a basis of $U\mf g$.
\end{theorem}
\begin{proof}[Sketch]
	Certainly $U\mf g$ is spanned by monomials, and one can use the relations $XY=YX+[X,Y]$ in order to rearrange the entries of the monomial. Thus, the given monomials span. Checking that they are linearly independent amounts to a long and technical argument.
\end{proof}

\subsection{Nilpotent and Solvable Lie Algebras}
\begin{definition}[ideal]
	Fix a Lie algebra $\mf g$. Then an \textit{ideal} $I\subseteq\mf g$ is a subspace such that $[X,I]\subseteq I$ for any $X\in\mf g$. Equivalently, $I$ is a $\mf g$-subrepresentation of $\mf g$.
\end{definition}
\begin{example}[center]
	Fix a Lie algebra $\mf g$. Then the \textit{center} is the subspace
	\[\mf z(\mf g)\coloneqq\{X\in\mf g:[X,\mf g]=0\}.\]
	The center is an ideal. It is kernel the adjoint representation $\op{ad}\colon\mf g\to\mf{gl}(\mf g)$.
\end{example}
\begin{example}[commutator]
	Fix a Lie algebra $\mf g$. Then the \textit{commutator} is the subspace $[\mf g,\mf g]$. By the Jacobi identity, $[\mf g,\mf g]$ is a Lie ideal.
\end{example}
\begin{remark}
	If $I$ is an ideal of a Lie algebra $\mf g$, the quotient space $\mf g/I$ is a Lie algebra, and there is a short exact sequence
	\[0\to I\to\mf g\to\mf g/I\to0\]
	of $\mf g$-modules.
\end{remark}
\begin{definition}
	A Lie algebra $\mf g$ is \textit{abelian} if and only if $[\mf g,\mf g]=0$.
\end{definition}
\begin{example}
	For any Lie algebra, the quotient $\mf g/[\mf g,\mf g]$ is abelian.
\end{example}
\begin{definition}[nilpotent]
	A Lie algebra $\mf g$ is \textit{nilpotent} if and only if there is a sequence
	\[\mf g=\mf a_0\supseteq\mf a_1\supseteq\cdots\supseteq\mf a_k=0\]
	such that $\mf a_{i+1}$ is an ideal of $\mf a_i$ with abelian quotient.
\end{definition}
\begin{remark}
	It is equivalent to check if the sequence $\mf a_{i+1}\coloneqq[\mf g,\mf a_i]$ eventually stabilizes at $0$.
\end{remark}
\begin{example}
	The subspace of strictly upper-triangular matrices of $\mf{gl}_n(\CC)$ is nilpotent.
\end{example}
\begin{definition}[solvable]
	A Lie algebra $\mf g$ is \textit{solvable} if and only if there is a sequence
	\[\mf g=\mf a_0\supseteq\mf a_1\supseteq\cdots\supseteq\mf a_k=0\]
	such that each $\mf a_i$ is an ideal of $\mf g$ such that $\mf g/\mf a_i$ is abelian.
\end{definition}
\begin{remark}
	It is equivalent to check if the sequence $\mf a_i\coloneqq[\mf a_{i+1},\mf a_{i+1}]$ eventually stabilizes at $0$.
\end{remark}
\begin{example}
	Any nilpotent Lie algebra is solvable.
\end{example}
\begin{example}
	The subspace of upper-triangular matrices of $\mf{gl}_n(\CC)$ is solvable but not nilpotent.
\end{example}
\begin{theorem}[Lie] \label{thm:lie}
	Let $\mf g$ be a solvable Lie algebra. For any representation $V$ of $\mf g$, there is a basis in which $\mf g$ acts by upper-triangular matrices.
\end{theorem}
\begin{proof}[Sketch]
	By induction on $\dim V$, it is enough to show that the action of $\mf g$ on $V$ admits a nonzero eigenvector. For this, we induct on $\dim\mf g$. By solvability, one may find a codimension-$1$ ideal $\mf g'\subseteq\mf g$; choose any $X\in\mf g\setminus\mf g'$. Now, induction tells us that $V$ admits a nonzero eigenvector $v$ for $\mf g'$. With some difficulty, can show that
	\[W\coloneqq\op{span}_\CC X^\NN v\]
	is a $\mf g$-subrepresentation of $V$, and it turns out that $\mf g'$ acts diagonally on all of $W$. Thus, any eigenvector of $X$ acting on $W$ will do.
\end{proof}
\begin{corollary}
	Let $\mf g$ be a Lie algebra. Then $\mf g$ is solvable if and only if $[\mf g,\mf g]$ is nilpotent.
\end{corollary}
\begin{proof}[Sketch]
	If $[\mf g,\mf g]$ is nilpotent, then one can show that $\mf g$ is solvable by building ideals from the short exact sequence
	\[0\to[\mf g,\mf g]\to\mf g\to\mf g/[\mf g,\mf g]\to0.\]
	Conversely, if $\mf g$ is solvable, then by \Cref{thm:lie}, one can find a basis of $\mf g$ on which the adjoint action of $\mf g$ on $\mf g$ is given by upper-triangular matrices. Thus, the action of $[\mf g,\mf g]$ is given by strictly upper-triangular matrices. This is enough to show $[\mf g,\mf g]$ is nilpotent.
\end{proof}
\begin{remark} \label{rem:nilp-to-strictly-upper-triangular}
	Similarly, if $\mf g$ is nilpotent, one can show that any representation $V$ admits a basis on which $\mf g$ acts by strictly upper-triangular matrices.
\end{remark}
\begin{theorem}[Engel] \label{thm:engel}
	Let $\mf g$ be a Lie algebra. Then $\mf g$ is nilpotent if and only if the operator $\op{ad}_X\colon\mf g\to\mf g$ is nilpotent for all $X\in\mf g$.
\end{theorem}
\begin{proof}[Sketch]
	If $\op{ad}_X$ is nilpotent for all $X\in\mf g$, then $[\mf g,[\mf g,[\mf g,\ldots,[\mf g,\mf g]]]]=0$ for a high enough nesting level, so $\mf g$ is nilpotent. Conversely, if $\mf g$ is nilpotent, then \Cref{rem:nilp-to-strictly-upper-triangular} grants a basis of $\mf g$ upon which $\mf g$ acts by strictly upper-triangular matrices, so $\op{ad}_X$ is nilpotent for all $X\in\mf g$.
\end{proof}

\subsection{Simple and Semisimple Lie Algebras}
\begin{definition}[simple]
	A Lie algebra $\mf g$ is \textit{simple} if and only if it is not abelian and admits no nonzero proper ideals.
\end{definition}
\begin{example}
	One can check by hand that $\mf{sl}_2(\CC)$ is a simple Lie algebra.
\end{example}
\begin{definition}[semisimple]
	A Lie algebra $\mf g$ is \textit{semismiple} if and only if it admits no nonzero solvable ideals.
\end{definition}
\begin{example}
	Any simple Lie algebra is semisimple. We will have many more examples shortly.
\end{example}
\begin{remark}
	It turns out that any semisimple Lie algebra $\mf g$ is a direct sum of simple ones. In short, one can decompose the adjoint action of $\mf g$ on $\mf g$ into a sum of irreducible subrepresentations, and each irreducible subrepresentation is a simple Lie subalgebra.
\end{remark}
\begin{remark}
	If $\mf g$ is semisimple, then $\mf z(\mf g)=0$, so the adjoint representation $\mf g\to\mf{gl}(\mf g)$ is injective.
\end{remark}
\begin{definition}[reductive]
	A Lie algebra $\mf g$ is \textit{reductive} if and only if $\mf g/\mf z(\mf g)$ is semisimple.
\end{definition}
\begin{example}
	Because $\mf z(\mf{gl}_2(\CC))$ consists of the scalar matrices, the Lie algebra $\mf{gl}_2(\CC)$ is reductive.
\end{example}
\begin{remark}
	A reductive Lie algebra $\mf g$ is semisimple if and only if $\mf z(\mf g)=0$.
\end{remark}
\begin{definition}[radical]
	Fix a Lie algebra $\mf g$. Then the \textit{radical} $\rad\mf g$ is the sum of all solvable ideals of $\mf g$. One can check that $\rad\mf g$ is a solvable ideal.
\end{definition}
\begin{example}
	For a Lie algebra $\mf g$, we see $\rad\mf g=0$ if and only if $\mf g$ is semisimple. For example, $\mf g/\rad\mf g$ is semisimple.
\end{example}
\begin{theorem}[Levi decomposition]
	Fix any Lie algebra $\mf g$. Then there is a semisimple Lie subalgebra $\mf g_{\mathrm{ss}}\subseteq\mf g$ and a direct sum decomposition
	\[\mf g=\rad\mf g\oplus\mf g_{\mathrm{ss}}.\]
\end{theorem}
\begin{proof}[Sketch]
	This is a very difficult result. The idea is to show that the short exact sequence
	\[0\to\rad\mf g\to\mf g\to\mf g_{\mathrm{ss}}\to0\]
	admits a splitting. Such an extension problem reduces to a calculation of some cohomology groups, which is more or less done by hand.
\end{proof}
\begin{proposition} \label{prop:get-reductive}
	Let $V$ be a representation of a Lie algebra $\mf g$. Suppose that the pairing $\langle-,-\rangle_V\colon\mf g\times\mf g\to\CC$ defined by
	\[\langle X,Y\rangle_V\coloneqq\tr(XY;V)\]
	is non-degenerate. Then $\mf g$ is reductive.
\end{proposition}
\begin{proof}[Sketch]
	It is enough to show that any $X\in[\mf g,\rad\mf g]$ lives in the kernel of $\langle-,-\rangle_V$ (and thus vanishes). By passing to some short exact sequences, we may assume that $V$ is irreducible. Using \Cref{thm:lie}, one can show that each $\rad(\mf g)$-eigenspace of $V$ is $\mf g$-invariant, so irreducibility implies that $\rad\mf g$ acts by scalars on $V$. A direct calculation now shows that $[\mf g,\rad\mf g]$ acts by zero on $V$.
\end{proof}
\begin{example}
	The standard representation $V=\CC^n$ of $\mf{gl}_n(\CC)$ satisfies the hypothesis of \Cref{prop:get-reductive}, so $\mf{gl}_n(\CC)$ is reductive. A similar argument holds for $\mf{so}_n(\CC)$.
\end{example}
\begin{theorem}[Cartan] \label{thm:cartan-criteria}
	Let $\mf g$ be a semisimple Lie algebra. Then the pairing
	\[\langle X,Y\rangle_{\mf g}\coloneqq\tr(XY;\mf g)\]
	is non-degenerate. This pairing is called the \textit{Killing form}.
\end{theorem}
\begin{proof}[Sketch]
	Please skip this sketch. Let $\mf h$ be the kernel of $\langle-,-\rangle_{\mf g}$, which is an ideal. We will show that $\mf h$ is a solvable, which forces $\mf h=0$. It is enough to show that $\op{ad}(\mf h)\subseteq\mf{gl}(\mf h)$ is solvable. Some linear algebra shows that the pairing $\langle X,Y\rangle\coloneqq\tr(XY;\mf h)$ on $\mf h$ vanishes. Thus, we are in the situation that we have a Lie algebra $\mf h'$ with a faithful representation $V$ for which
	\[\tr(XY;V)=0\]
	for each $X,Y\in\mf h'$. One can use the Jordan decomposition to show that any element of $[\mf h',\mf h']$ is nilpotent, so $[\mf h',\mf h']$ is nilpotent, so $\mf h'$ is solvable.
\end{proof}

\subsection{The Root Decomposition}
\begin{definition}
	Fix a Lie algebra $\mf g$. Then $X\in\mf g$ is \textit{nilpotent} or \textit{semisimple} if and only if $\op{ad}_X\colon\mf g\to\mf g$ is nilpotent or semisimple, respecively.
\end{definition}
\begin{proposition}[Jordan decomposition]
	Fix a semisimple Lie algebra $\mf g$. For any $X\in\mf g$, there is a unique semisimple $X_s\in\mf g$ and $X_n\in\mf g$ such that $[X_s,X_n]=0$ and $X=X_s+X_n$.
\end{proposition}
\begin{proof}[Sketch]
	The usual Jordan decomposition produces a unique pair $(X_s',X_n')\in\mf{gl}(\mf g)^2$ for which $X_s'$ is semisimple, $X_n'$ is nilpotent, $[X_s',X_n']=0$, and $\mathrm{ad}_X=X_s'+X_n'$. Thus, it remains to check that $X_s'$ and $X_n'$ are in the image of $\op{ad}\colon\mf g\to\mf{gl}(\mf g)$. In short, the Jordan decomposition shows that $X_s'$ and $X_n'$ are polynomials in $\op{ad}_X$, so they are derivations of $\mf g$, but all derivations of $\mf g$ come from $\mf g$.
\end{proof}
\begin{definition}[Cartan]
	Fix a semisimple Lie algebra $\mf g$. Then a \textit{torus} is an abelian subalgebra $\mf h$ of $\mf g$ consisting of semisimple elements. A maximal torus is a \textit{Cartan subalgebra}.
\end{definition}
\begin{example}
	The subalgebra of diagonal matrices of $\mf{sl}_n(\CC)$ is a Cartan subalgebra.
\end{example}
\begin{remark}
	Let $\mf h$ be a Cartan subalgebra. It turns out that
	\[\mf h=\{X\in\mf h:[X,\mf h]=0\}.\]
	In short, use the Jordan decomposition and the maximality of $\mf h$.
\end{remark}
\begin{remark}
	It turns out that any two Cartan subalgebras of $\mf g$ are conjugate under $\op{Aut}\mf g$. This is fairly technical, and we will only use it to show that certain constructions are independent of the choice of $\mf h$.
\end{remark}
\begin{definition}[rank]
	Fix a semisimple Lie algebra $\mf g$. Then the \textit{rank} of $\mf g$ is the dimension of a Cartan subalgebra of $\mf g$.
\end{definition}
\begin{example}
	The rank of $\mf{sl}_n(\CC)$ is $n-1$.
\end{example}
\begin{definition}
	Fix a semisimple Lie algebra $\mf g$ and a Cartan subalgebra $\mf h\subseteq\mf g$. Then the \textit{root system} $\Phi\subseteq\mf h^*$ is the set of nonzero eigenspaces of $\mf h$ acting on $\mf g$.
\end{definition}
\begin{proposition}
	Fix a semisimple Lie algebra $\mf g$ and a Cartan subalgebra $\mf h\subseteq\mf g$. Then the \textit{root decomposition} is the eigenspace decomposition
	\[\mf g=\mf h\oplus\bigoplus_{\alpha\in\Phi}\mf g_\alpha,\]
	where $\mf g_\alpha$ is the eigenspace of $\alpha\in\Phi$.
\end{proposition}
\begin{proof}[Sketch]
	Diagonalize the representation $\mf g$ of $\mf h$.
\end{proof}
\begin{remark}
	For $\alpha,\beta\in\Phi$, we have $[\mf g_\alpha,\mf g_\beta]\subseteq\mf g_{\alpha+\beta}$.
\end{remark}
\begin{remark}
	The subspace $\Phi$ spans $\mf h^*$. Indeed, any $h\in\mf h$ orthogonal to all $\alpha\in\Phi$ then has ${\op{ad}_h}=0$, so $h\in\mf z(\mf g)$, so $h=0$.
\end{remark}
\begin{notation}
	Fix a semisimple Lie algebra $\mf g$ and a Cartan subalgebra $\mf h\subseteq\mf g$ and root system $\Phi\subseteq\mf h^*$. For any $\alpha\in\Phi$, the subspaces $\mf g_\alpha$ and $\mf g_{-\alpha}$ generate a subalgebra of $\mf g$ isomorphic to $\mf{sl}_2(\CC)$. (This is a little tricky to show.) We denote this subalgebra by $\mf{sl}_2(\CC)_\alpha$.
\end{notation}
\begin{proposition}
	Fix a semisimple Lie algebra $\mf g$ and a Cartan subalgebra $\mf h\subseteq\mf g$ and root system $\Phi\subseteq\mf h^*$. Let $(-,-)$ be the Killing form on $\mf g$.
	\begin{listalph}
		\item For each $\alpha\in\Phi$, $\dim\mf g_\alpha=1$.
		\item For each $\alpha,\beta\in\Phi$, $2(\alpha,\beta)/(\alpha,\alpha)\in\ZZ$.
		\item For each $\alpha\in\Phi$, define the reflection $s_\alpha\colon\mf h^*\to\mf h^*$ defined by $s_\alpha(\lambda)\coloneqq\lambda-2(\alpha,\lambda)/(\alpha,\alpha)\cdot\alpha$. Then $s_\alpha$ preserves $\Phi$.
		\item If $\alpha\in\Phi$ and $k\alpha\in\Phi$ for some $k\in\CC$, then $k\in\{+1,-1\}$.
	\end{listalph}
\end{proposition}
\begin{proof}[Sketch]
	Let $(-,-)$ be the Killing form on $\mf h^*$, which is non-degenerate. Define $h_\alpha\coloneqq2H_\alpha/(\alpha,\alpha)$, where $H_\alpha$ is the element dual to $\alpha\in\mf h^*$. Then one can show that
	\[\CC h_\alpha\oplus\bigoplus_{k\in\ZZ\setminus\{0\}}\mf g_{k\alpha}\]
	is an irreducible representation of $\mf{sl}_2(\CC)_\alpha$. Each part now follows from understanding the irredcible representations of $\mf{sl}_2(\CC)$.
\end{proof}

% \subsection{Root Systems}
% \begin{definition}
	
% \end{definition}

\subsection{Representation Theory of Semisimple Lie Algebras}
\begin{theorem}
	Let $\mf g$ be a semisimple Lie algebra. Then all representations of $\mf g$ are completely reducible.
\end{theorem}
\begin{proof}[Sketch]
	There are many proofs of this theorem. The main idea is to show that any short exact sequence
	\[0\to V'\to V\to V''\to0\]
	of representations of $\mf g$ splits as a direct sum $V=V'\oplus V''$. The most algebraic proof proceeds by finding ``Casimir'' elements in the center of $U\mf g$. By \Cref{lem:schur}, they act by scalars on irreducible representations. One can force the splitting if granted a Casimir elements which acts on $V'$ and $V''$ by different scalars.
\end{proof}

\end{document}